\section{Global convergence and discretization}
\label{sec:global}

We study the Gaussian natural gradient flow where the target posterior distribution
$\rho_\post$ is logconcave, and then propose practical discretization algorithms.

\subsection{Setup and the flow}

\begin{assumption}\label{assump:logconcave-smooth}
    We assume that $\rho_\post$ is $\alpha$-strongly logconcave and $\beta$-smooth, i.e.
    \begin{equation*}
        \alpha \Id \preceq -\nabla_\theta \nabla_\theta \log \rho_\post (\theta) \preceq \beta \Id,
        \qquad \theta\in\R^{N_\theta},
    \end{equation*}
    which is equivalent to
    \begin{equation*}
        \alpha \Id \preceq \nabla^2 V(\theta) \preceq \beta \Id .
    \end{equation*}
    We write $\kappa := \beta/\alpha \geq 1$ for the condition number.
\end{assumption}

A third of our results, all of them stated where they are used, need in addition that the
Hessian varies slowly.

\begin{assumption}[Hessian-Lipschitzness]
\label{assump:hess-lip}
There is $\LH<\infty$ with $\norm{\nabla^2V(x)-\nabla^2V(y)}_\op\leq\LH\norm{x-y}$ for all
$x,y\in\R^{N_\theta}$. We write the dimensionless parameter $\barLH:=\LH/\alpha^{3/2}$; a Gaussian
target has $\LH=\barLH=0$.
\end{assumption}

For a Gaussian parameter $a=(m,C)\in\Aspace$ we abbreviate the averaged score and the
averaged curvature,
\begin{equation}\label{eq:gA}
    g(a) := \E_{\rho_a}[\nabla_\theta\log\rho_\post(\theta)]
          = -\E_{\rho_a}[\nabla V(\theta)],
    \qquad
    A(a) := \E_{\rho_a}[\nabla^2V(\theta)],
\end{equation}
and the precision residual $R(a) := C^{-1}-A(a)$. Under
\Cref{assump:logconcave-smooth} we have $\alpha\Id\preceq A(a)\preceq\beta\Id$ for every
$a\in\Aspace$, and at the minimizer the first-order conditions read
\begin{equation}\label{eq:first-order}
    g(a_\star)=0,
    \qquad
    A(a_\star)=C_\star^{-1},
    \qquad\text{that is}\qquad
    R(a_\star)=0 .
\end{equation}
In this notation the Gaussian natural gradient flow \eqref{ODEs:NG} reads
\begin{equation}\label{eq:NG-compact}
    \frac{\dd m_t}{\dd t}=C_tg(a_t),
    \qquad
    \frac{\dd C_t}{\dd t}=C_t-C_tA(a_t)C_t .
\end{equation}
We also use the energy splitting $\calE = \calE_1+\calE_2$ with
\begin{equation}\label{eq:E-split}
    \calE_1(a)=\int\rho_a\log\rho_a
    =-\tfrac12\log\det C+\mathrm{const.},
    \qquad
    \calE_2(a)=-\E_{\rho_a}[\log\rho_\post]
    =\E_{\rho_a}[V]+\mathrm{const.}
\end{equation}

The Fisher--Rao metric induced on the Gaussian family, at $a=(m,C)$ and for tangent vectors
$(u,U),(v,W)\in\R^{N_\theta}\times\Sym(N_\theta)$, is
\begin{equation}\label{eq:FR-metric}
    \inner{(u,U)}{(v,W)}_a
    :=
    u^\top C^{-1}v+\tfrac12\Tr\bigl(C^{-1}UC^{-1}W\bigr),
    \qquad
    \norm{(u,U)}_a^2
    =
    u^\top C^{-1}u+\tfrac12\norm{C^{-1/2}UC^{-1/2}}_\Frob^2,
\end{equation}
which is the Fisher information inner product of the Gaussian family. We first identify the
flow \eqref{ODEs:NG} as the gradient flow of $\calE$ in this metric. Writing
$\theta=m+C^{1/2}Z$ with $Z\sim\N(0,\Id)$ and differentiating under the integral sign, which
is justified by dominated convergence since $\norm{\nabla^2V}_\op\leq\beta$ forces
$\norm{\nabla V(\theta)}\leq\norm{\nabla V(0)}+\beta\norm\theta$, we obtain
$\partial_m\calE_2=-g(a)$ and, by Price's theorem,
$\partial_C\calE_2=\tfrac12A(a)$. Together with
$\partial_C\calE_1=-\tfrac12C^{-1}$ this gives the differential
\begin{equation}\label{eq:differential}
    D\calE(a)[(u,U)]
    =
    -g(a)^\top u+\tfrac12\Tr\bigl((A(a)-C^{-1})U\bigr).
\end{equation}
Matching \eqref{eq:differential} against \eqref{eq:FR-metric} block by block, the Riemannian
gradient is
\begin{equation}\label{eq:FR-gradient}
    \grad\calE(a)=\bigl(-Cg(a),\;CA(a)C-C\bigr),
\end{equation}
so that $\dot a_t=-\grad\calE(a_t)$ is exactly \eqref{ODEs:NG}. Its squared norm is
\begin{equation}\label{eq:grad-norm}
    \norm{\grad\calE(a)}_a^2
    =
    \Gradnorm(a)^2
    :=
    g(a)^\top Cg(a)+\tfrac12\norm{C^{1/2}R(a)C^{1/2}}_\Frob^2,
\end{equation}
because $C^{-1}(CA(a)C-C)=A(a)C-\Id$ and cyclicity of the trace turn the covariance block
into $\Tr(CR(a)CR(a))$. The chain rule then gives the exact dissipation identity
\begin{equation}\label{eq:dissipation}
    \frac{\dd}{\dd t}\DeltaE(a_t)
    =
    D\calE(a_t)[-\grad\calE(a_t)]
    =
    -\Gradnorm(a_t)^2 .
\end{equation}

Two inequalities convert \eqref{eq:dissipation} into a rate. The first is the Wasserstein
Polyak--{\L}ojasiewicz inequality induced by the $\alpha$-strong convexity of $V$ and
restricted to the Gaussian family \citep{lambert2022variational}: for every $a\in\Aspace$,
\begin{equation}\label{eq:W2-PL}
    \norm{\grad^{W_2}\calE(a)}_{W_2}^2
    =
    \norm{g(a)}^2+\Tr\bigl(R(a)CR(a)\bigr)
    \geq
    2\alpha\,\DeltaE(a).
\end{equation}
The second compares the two gradient norms. If $C\succeq\ell\Id$, then
$g(a)^\top Cg(a)\geq\ell\norm{g(a)}^2\geq\tfrac\ell2\norm{g(a)}^2$, while
$R(a)CR(a)\succeq0$ gives
$\tfrac12\Tr(CR(a)CR(a))=\tfrac12\Tr(C[R(a)CR(a)])\geq\tfrac\ell2\Tr(R(a)CR(a))$; adding the
two estimates and inserting \eqref{eq:W2-PL},
\begin{equation}\label{eq:FR-W2}
    \Gradnorm(a)^2
    \geq
    \frac\ell2\norm{\grad^{W_2}\calE(a)}_{W_2}^2
    \geq
    \alpha\ell\,\DeltaE(a),
\end{equation}
and the same computation with $\norm C_\op$ in place of $\ell$ gives the reverse comparison
$\Gradnorm(a)^2\leq\norm C_\op\norm{\grad^{W_2}\calE(a)}_{W_2}^2$. We combine \eqref{eq:dissipation} and \eqref{eq:FR-W2} with a lower bound on the
covariance.

\subsection{Continuous-time global convergence}

The covariance lower bound is cleanest in the precision variable $P_t:=C_t^{-1}$, along which
\eqref{ODEs:NG} is affine: differentiating $P_t=C_t^{-1}$ and using \eqref{eq:NG-compact},
\begin{equation}\label{eq:precision-duhamel}
    \dot P_t=A(a_t)-P_t,
    \qquad
    P_t=e^{-t}P_0+\int_0^te^{-(t-s)}A(a_s)\,\dd s .
\end{equation}

The two-sided covariance envelope along the flow, \Cref{lem:cov-cont}, is proved in
\Cref{app:det-proofs}.

The rate produced by \eqref{eq:FR-W2} and \Cref{lem:cov-cont} (\Cref{app:det-proofs}) still carries the affine
anisotropy of the target, which the flow itself does not see. We therefore whiten by the
optimizing Gaussian covariance. Let $a_\star=(m_\star,C_\star)$ be the optimizing Gaussian and
set
\begin{equation}\label{eq:whitening}
    \widehat m:=C_\star^{-1/2}(m-m_\star),
    \qquad
    \widehat C:=C_\star^{-1/2}CC_\star^{-1/2},
    \qquad
    \widehat V(z):=V(m_\star+C_\star^{1/2}z),
\end{equation}
so that the optimizer becomes $(\widehat m_\star,\widehat C_\star)=(0,\Id)$, and define the
optimal whitened curvature constants and the whitened initial covariance floor,
\begin{equation}\label{eq:star-curvature}
    \alphastar:=\inf_{z}\lambda_{\min}(\nabla^2\widehat V(z)),
    \qquad
    \betastar:=\sup_{z}\lambda_{\max}(\nabla^2\widehat V(z)),
    \qquad
    \kappastar:=\frac{\betastar}{\alphastar},
\end{equation}
\begin{equation}\label{eq:lamzstar}
    \lamzstar:=\lambda_{\min}(\widehat C_0)
    =\lambda_{\min}\bigl(C_\star^{-1/2}C_0C_\star^{-1/2}\bigr).
\end{equation}
Writing $S=C_\star^{1/2}$ we have $\nabla\widehat V(z)=S\nabla V(m_\star+Sz)$ and
$\nabla^2\widehat V(z)=S\nabla^2V(m_\star+Sz)S$, hence $\widehat g=Sg$ and
$\widehat A=SAS$, so the whitened parameters solve the same Fisher--Rao system
$\dot{\widehat m}=\widehat C\widehat g$,
$\dot{\widehat C}=\widehat C-\widehat C\widehat A\widehat C$. The first-order conditions
\eqref{eq:first-order} become
$\E_{Z\sim\N(0,\Id)}[\nabla\widehat V(Z)]=0$ and
$\E_{Z\sim\N(0,\Id)}[\nabla^2\widehat V(Z)]=\Id$, so that
$\alphastar\leq1\leq\betastar$. Under an invertible affine change $y=Bx+b$ the transformed
optimizer covariance is $BC_\star B^\top$, and
$O:=(BC_\star B^\top)^{-1/2}BC_\star^{1/2}$ is orthogonal; whitening before and after the
change therefore differs by an orthogonal map only, so $\alphastar$, $\betastar$,
$\kappastar$ and the spectrum of $\widehat C_0$, hence also the product
$\betastar\lamzstar$, are affine invariants. Finally, $C_\star^{-1}=A(a_\star)$ and
$\alpha\Id\preceq A(a_\star)\preceq\beta\Id$ give
$\beta^{-1}\Id\preceq C_\star\preceq\alpha^{-1}\Id$, so
$\kappa^{-1}\Id\preceq\nabla^2\widehat V(z)\preceq\kappa\Id$ and
\begin{equation}\label{eq:kappa-star-kappa}
    1\leq\kappastar\leq\kappa^2 .
\end{equation}

\begin{theorem}[Affine-invariant continuous-time global convergence]\label{thm:glob-cont}
    Under \Cref{assump:logconcave-smooth}, the Gaussian natural gradient flow
    \eqref{ODEs:NG} satisfies, for every $t\geq0$,
    \begin{equation*}
        \DeltaE(a_t)
        \leq
        \DeltaE_0\bigl[1+\betastar\lamzstar(e^t-1)\bigr]^{-1/\kappastar}.
    \end{equation*}
    Consequently, for every $\delta\in(0,\DeltaE_0)$ the hitting time
    $T^{\glob}(\delta):=\inf\{t\geq0:\DeltaE(a_t)\leq\delta\}$ obeys
    \begin{equation*}
        T^{\glob}(\delta)
        \leq
        \log\left(1+\frac{(\DeltaE_0/\delta)^{\kappastar}-1}{\betastar\lamzstar}\right)
        \leq
        \logp\frac1{\betastar\lamzstar}
        +\kappastar\log\frac{\DeltaE_0}{\delta}.
    \end{equation*}
    The original-coordinate estimate remains valid as well, so we may take the minimum of the
    two, replacing $(\betastar,\lamzstar,\kappastar)$ by $(\beta,\lambda_0,\kappa)$.
\end{theorem}

\begin{proof}
    The energy and the Fisher--Rao dissipation are invariant under the whitening
    \eqref{eq:whitening}, so we work in whitened coordinates and omit hats. Applying
    \Cref{lem:cov-cont} with the whitened constants gives
    $C_t\succeq\ell_\star(t)\Id$ with
    $\ell_\star(t)=\lamzstar e^t/(1+\betastar\lamzstar(e^t-1))$, and \eqref{eq:FR-W2}
    gives $\Gradnorm(a_t)^2\geq\alphastar\ell_\star(t)\DeltaE(a_t)$. With
    \eqref{eq:dissipation} this yields
    $\dot\DeltaE_t\leq-\alphastar\ell_\star(t)\DeltaE_t$, and by Gronwall's inequality
    \begin{equation*}
        \DeltaE_t
        \leq
        \DeltaE_0\exp\left(-\alphastar\int_0^t\ell_\star(s)\,\dd s\right).
    \end{equation*}
    The scalar integral is exact,
    $\int_0^t\ell_\star(s)\,\dd s=\betastar^{-1}\log[1+\betastar\lamzstar(e^t-1)]$, and
    $\alphastar/\betastar=1/\kappastar$, which gives the decay estimate. Solving the equality
    in it gives the exact hitting time. For the simplified form, set
    $R=(\DeltaE_0/\delta)^{\kappastar}\geq1$ and $b=\betastar\lamzstar$: if $b\leq1$ then
    $1+(R-1)/b\leq R/b$, and if $b\geq1$ then $1+(R-1)/b\leq R$; taking logarithms in the two
    cases gives the bound.
\end{proof}

The intrinsic condition number removes anisotropy that can be eliminated by an affine
reparametrization. It does not, by itself, encode the initialization size: the factor
$\log(\DeltaE_0/\delta)$, or an equivalent radius term, is indispensable, since if
arbitrary initial means are allowed then the worst-case time to reach a fixed energy level
is infinite even for a quadratic target. The comparison in \eqref{eq:kappa-star-kappa} can
be very loose in the useful direction: for an anisotropic quadratic target
$\kappastar=1$ while $\kappa$ is arbitrarily large.

\begin{remark}
\label{rem:FR-init-cov}
From Corollary 3 of \citet{lambert2022variational}, the convergence rate of Gaussian
approximate Wasserstein gradient flow does not depend on the initial covariance, while in
\Cref{thm:glob-cont} the convergence rate of the Gaussian approximate natural gradient flow
depends on the spectral bound of the initial covariance through \Cref{lem:cov-cont}.
We show by example that the dependence on the lower spectral bound of the initial covariance is
not a proof artifact, instead, it is determined by the properties of the Fisher--Rao gradient
flow. For the one-dimensional Gaussian target
\[
    \rho_{\post}=\N(0,1),
    \qquad
    V(\theta)=\frac12\theta^2,
\]
with $a=(m,c)$, the Gaussian variational objective is
\[
    \calE(m,c)-\calE(a_\star)
    =
    \frac12\bigl(m^2+c-1-\log c\bigr),
    \qquad
    a_\star=(0,1).
\]
Since
\[
    \nabla\log\rho_{\post}(\theta)=-\theta,
    \qquad
    \nabla^2\log\rho_{\post}(\theta)=-1,
\]
the Gaussian natural gradient flow becomes
\[
    \dot m_t=-c_tm_t,
    \qquad
    \dot c_t=c_t(1-c_t).
\]
Let $c_0=\varepsilon\in(0,1)$. Then the covariance equation is exactly solvable:
\[
    c_t=\frac{\varepsilon e^t}{1-\varepsilon+\varepsilon e^t}.
\]
Consequently,
\[
    m_t
    =
    m_0\exp\left(-\int_0^tc_s\,\dd s\right)
    =
    \frac{m_0}{1-\varepsilon+\varepsilon e^t}.
\]
Therefore, for every fixed $t>0$,
\[
    \lim_{m_0\to\infty}
    \frac{\calE(m_t,c_t)-\calE(a_\star)}{\calE(m_0,c_0)-\calE(a_\star)}
    =
    \frac1{(1-\varepsilon+\varepsilon e^t)^2}.
\]
Taking $\varepsilon\downarrow0$ gives
\[
    \lim_{\varepsilon\downarrow0}\lim_{m_0\to\infty}
    \frac{\calE(m_t,c_t)-\calE(a_\star)}{\calE(m_0,c_0)-\calE(a_\star)}
    =1 .
\]
Hence no estimate of the form
\[
    \calE(a_t)-\calE(a_\star)
    \leq
    e^{-rt}\bigl(\calE(a_0)-\calE(a_\star)\bigr),
    \qquad r>0,
\]
with $r$ depending only on the log-concavity and smoothness constants, can hold uniformly over
all initial covariances. In this example $\alpha=\beta=1$, and the rate in the preceding result
depends on
\[
    \lambda_{\min}=\min\{\varepsilon,1\}=\varepsilon,
\]
which correctly degenerates as $\varepsilon\downarrow0$.

This should be contrasted with the Gaussian Bures--Wasserstein gradient flow for the same
target. In one dimension it satisfies
\[
    \dot{\widetilde m}_t=-\widetilde m_t,
    \qquad
    \dot{\widetilde c}_t=2(1-\widetilde c_t),
\]
and hence
\[
    \widetilde m_t=\widetilde m_0e^{-t},
    \qquad
    \widetilde c_t=1+(\widetilde c_0-1)e^{-2t}.
\]
Thus the mean component contracts immediately at the target curvature rate, independently of
the initial covariance. This illustrates the structural difference between the Fisher--Rao and
Bures--Wasserstein geometries: in the Fisher--Rao Gaussian flow, small covariance preconditions
and slows down the mean dynamics, while the Bures--Wasserstein flow does not contain this
covariance prefactor in the mean equation. The same contrast is what makes a short transport
warm start remove the covariance burn-in term altogether, as we show in
\Cref{subsec:burnin}.
\end{remark}

The linear dependence on $\kappastar$ in \Cref{thm:glob-cont} cannot be improved, up to
logarithmic initialization and localization factors, under the general spectral initialization
class.

\begin{theorem}[Continuous-time logarithmic-spiral lower bound]\label{thm:sharp-cont}
    For every $K\geq16$ there is an even potential $V_K\in C^3(\R^2)$, given explicitly in
    \eqref{eq:spiral-potential}, and a scale $r_0$ with
    $r_0\sqrt\alpha\geq10^{25}K^{3/2}\sqrt{\log(eK)}$, such that
    $\alpha\Id\preceq\nabla^2V_K\preceq\beta\Id$ with $\kappa=2K+1$ and
    \begin{equation*}
        K\leq\kappastar(V_K)\leq1.03K .
    \end{equation*}
    There is an initial mean $m_0$ with $\norm{m_0}=e^2r_0$ and an anisotropic initial
    covariance $C_0$, admissible under the general spectral condition
    $\beta^{-1}\Id\preceq C_0\preceq\alpha^{-1}\Id$, such that for every local level
    $\delta_{\loc}<\tfrac\alpha2e^{-0.08}\norm{m_0}^2$ the exact Gaussian Fisher--Rao flow
    obeys
    \begin{equation*}
        T_{\loc}:=\inf\{t\geq0:\DeltaE(a_t)\leq\delta_{\loc}\}
        \geq\frac K{100}
        \geq\frac{\kappastar(V_K)}{103}.
    \end{equation*}
    Moreover $\betastar\lamzstar\geq(1-\eta_K)/(1+\eta_K)$ with
    $\eta_K:=Ke^{-\alpha r_0^2/2}\leq10^{-3}$, so the covariance-initialization term
    $\logp(1/(\betastar\lamzstar))$ in \Cref{thm:glob-cont} is smaller than $0.003$ for this
    family.
\end{theorem}

The proof is given in \Cref{app:spiral}.

A $\Theta(\kappastar)$ packaging of this bound, \Cref{cor:sharp-cont}, is stated in
\Cref{app:spiral}.

The matching construction uses an admissible anisotropic initial covariance. It does not
establish an $\Omega(\kappastar)$ lower bound under the stricter initialization
$C_0=\beta^{-1}\Id$. It also allows the initial radius to grow with $K$, and the logarithm
of that radius is part of the input size charged by \Cref{thm:glob-cont}. A
Hessian-Lipschitz version of the construction is available for every prescribed
$\LH>0$ by enlarging $r_0$, and is recorded in \Cref{app:spiral}.

\subsection{Discretization with Riemannian retraction}

We decompose the objective as $\calE=\calE_1+\calE_2$ as in \eqref{eq:E-split}, and define
the squared base-point penalty at the current iterate $a_n=(m_n,C_n)$ as
\begin{equation}\label{eq:Riem-dist}
    d(a,a_n)^2
    =
    (m-m_n)^\top C_n^{-1}(m-m_n)+\frac12\norm Z_\Frob^2,
    \qquad
    Z=\log\bigl(C_n^{-1/2}CC_n^{-1/2}\bigr),
\end{equation}
built blockwise from the metric \eqref{eq:FR-metric} at the base point $a_n$: the mean block
freezes the covariance at $C_n$, and the covariance block is the affine-invariant distance on
$\SPD(N_\theta)$. It is not the geodesic distance of \eqref{eq:FR-metric}, whose geodesics
couple the mean and the covariance (\Cref{rem:Riem-retr}). As the entropy term $\calE_1$ is
often non-smooth, we apply the proximal algorithm \citep{parikh2014proximal} for optimization.
In the forward step,
\begin{equation*}
    a_{n+1/2}
    =
    \argmin_{a\in\Aspace}
    \left\{
        \calE_2(a_n)+\inner{\grad\calE_2(a_n)}{\Retr_{a_n}^{-1}(a)}+\frac1{2\Delta t}d(a,a_n)^2
    \right\},
\end{equation*}
where $\Retr_{a_n}^{-1}$ is the inverse of the retraction $\Retr_{a_n}$ defined in
\Cref{prop:upd-map-Riem} below, which gives
\begin{equation*}
    m_{n+1/2}=m_n+\Delta t\,C_ng(a_n),
    \qquad
    C_{n+1/2}=C_n^{1/2}\exp\bigl(-\Delta t\,C_n^{1/2}A(a_n)C_n^{1/2}\bigr)C_n^{1/2}.
\end{equation*}
The backward step reads
\begin{equation*}
    a_{n+1}
    =
    \argmin_{a\in\Aspace}
    \left\{\calE_1(a)+\frac1{2\Delta t}d(a,a_{n+1/2})^2\right\},
\end{equation*}
which gives
\begin{equation*}
    m_{n+1}=m_{n+1/2},
    \qquad
    C_{n+1}=e^{\Delta t}C_{n+1/2}.
\end{equation*}
These collectively give the update through Riemannian geometry:
\begin{equation}\label{upd:Riem}
    \boxed{
    \begin{aligned}
        m_{n+1} &= m_n+\Delta t\,C_ng(a_n), \\
        C_{n+1} &= C_n^{1/2}
        \exp\bigl(\Delta t(\Id-C_n^{1/2}A(a_n)C_n^{1/2})\bigr)C_n^{1/2}.
    \end{aligned}}
\end{equation}

We proceed to show the convergence theory of \eqref{upd:Riem} under
\Cref{assump:logconcave-smooth}. We first show that the update for each step is exactly a
gradient descent step along the exponential retraction of the Gaussian manifold.

\begin{proposition}\label{prop:upd-map-Riem}
    For $a=(m,C)\in\Aspace$ and $(u,U)\in T_a\Aspace$ define
    \begin{equation*}
        \Retr_{(m,C)}(u,U)
        :=
        \left(m+u,\;C^{1/2}\exp\bigl(C^{-1/2}UC^{-1/2}\bigr)C^{1/2}\right).
    \end{equation*}
    Then the discrete algorithm with Riemannian retraction \eqref{upd:Riem} gives the one-step
    update
    \begin{equation*}
        a_{n+1}=\Retr_{a_n}\bigl(-\Delta t\grad\calE(a_n)\bigr).
    \end{equation*}
\end{proposition}

\begin{proof}
    By \eqref{eq:FR-gradient}, in the mean direction
    \begin{equation*}
        m_{n+1}=m_n+\Delta t\,C_ng(a_n)=m_n-\Delta t\grad_m\calE(a_n).
    \end{equation*}
    In the covariance direction,
    \begin{align*}
        -\Delta t\grad_C\calE(a_n)
        &=\Delta t\,C_n-\Delta t\,C_nA(a_n)C_n\\
        &=C_n^{1/2}\bigl(\Delta t\,\Id-\Delta t\,C_n^{1/2}A(a_n)C_n^{1/2}\bigr)C_n^{1/2},
    \end{align*}
    so that $C_n^{-1/2}(-\Delta t\grad_C\calE(a_n))C_n^{-1/2}
    =\Delta t(\Id-C_n^{1/2}A(a_n)C_n^{1/2})$, which is exactly the exponent appearing in the
    covariance component of \eqref{upd:Riem}. This completes the proof.
\end{proof}

\begin{remark}\label{rem:Riem-retr}
    The covariance block of $\Retr$ is the exponential map of the affine-invariant metric on
    $\SPD(N_\theta)$, so the covariance update follows a genuine geodesic of the covariance
    factor. The joint mean--covariance map is nevertheless a retraction and not the exponential
    map of the Fisher--Rao metric on $\Aspace$, whose geodesics couple the mean and the
    covariance. All estimates below are proved for the retraction directly, so no geodesic
    smoothness statement is used.
\end{remark}

We then prove the spectral bounds of the covariance along \eqref{upd:Riem}.

\begin{lemma}[Riemannian covariance bootstrap]\label{lem:cov-Riem}
    Under \Cref{assump:logconcave-smooth}, if
    $\lambda_{0,\min}\Id\preceq C_0\preceq\lambda_{0,\max}\Id$ and
    $0<\Delta t\leq\min\{1,1/(\beta\lambda_{\max})\}$, then along \eqref{upd:Riem} one has
    $C_n\preceq\lambda_{\max}\Id$ for all $n\geq0$, and $C_n\succeq L_n\Id$ where
    \begin{equation*}
        L_0:=\min\left\{\lambda_0,\frac1{2\beta}\right\},
        \qquad
        L_{n+1}:=\frac{e^{\Delta t}L_n}{1+(e-1)\Delta t\,\beta L_n}.
    \end{equation*}
    With $y_n:=(\beta L_n)^{-1}$, $y_0=1/(\beta L_0)$ and
    $y_\infty:=(e-1)\Delta t/(e^{\Delta t}-1)\in[1,e-1)$, the envelope solves the affine
    recursion $y_n=y_\infty+(y_0-y_\infty)e^{-n\Delta t}$. Hence
    $C_n\succeq\tfrac1{2\beta}\Id$ for all
    \begin{equation*}
        n\geq
        N_{\cov,-}^{\Riem}
        :=
        \left\lceil\frac1{\Delta t}\logp\frac{y_0-y_\infty}{2-y_\infty}\right\rceil
        =
        O\!\left(\frac1{\Delta t}\logp\frac1{\beta\lambda_0}\right).
    \end{equation*}
    Moreover, with $z_0:=1/(\alpha\lambda_{0,\max})$ and
    $z_\infty:=\Delta t/(e^{\Delta t}-1)\in[\tfrac1{e-1},1)$, one has
    $z_n:=\lambda_{\min}(C_n^{-1})/\alpha\geq z_\infty+(z_0-z_\infty)e^{-n\Delta t}$, hence
    $C_n\preceq\tfrac2\alpha\Id$ for all $n\geq N_{\cov,+}^{\Riem}$, where
    $N_{\cov,+}^{\Riem}:=0$ if $z_0\geq\tfrac12$ and
    $N_{\cov,+}^{\Riem}:=\lceil\Delta t^{-1}\log\frac{z_\infty-z_0}{z_\infty-1/2}\rceil
    =O(\Delta t^{-1}\logp(\alpha\lambda_{0,\max}))$ otherwise.
\end{lemma}

\begin{proof}
    Write $B_n:=C_n^{1/2}A(a_n)C_n^{1/2}$, so that
    $C_{n+1}^{-1}=C_n^{-1/2}\exp(\Delta t(B_n-\Id))C_n^{-1/2}$. Using $e^X\succeq\Id+X$ for
    symmetric $X$ and $C_n^{-1/2}B_nC_n^{-1/2}=A(a_n)$,
    \begin{equation*}
        C_{n+1}^{-1}\succeq(1-\Delta t)C_n^{-1}+\Delta t\,A(a_n).
    \end{equation*}
    If $C_n\preceq\lambda_{\max}\Id$ then $C_n^{-1}\succeq\lambda_{\max}^{-1}\Id$ and
    $A(a_n)\succeq\alpha\Id\succeq\lambda_{\max}^{-1}\Id$, and since $\Delta t\leq1$ this gives
    $C_{n+1}^{-1}\succeq\lambda_{\max}^{-1}\Id$, which proves the upper bound by induction.
    For the lower envelope, $0\preceq\Delta t\,B_n\preceq\Id$ because
    $\Delta t\leq1/(\beta\lambda_{\max})$, and the scalar inequality $e^x\leq1+(e-1)x$ on
    $[0,1]$ gives $\exp(\Delta t\,B_n)\preceq\Id+(e-1)\Delta t\,B_n$ by functional calculus,
    hence
    \begin{equation*}
        C_{n+1}^{-1}
        =e^{-\Delta t}C_n^{-1/2}\exp(\Delta t\,B_n)C_n^{-1/2}
        \preceq e^{-\Delta t}\bigl[C_n^{-1}+(e-1)\Delta t\,A(a_n)\bigr]
        \preceq e^{-\Delta t}\bigl(L_n^{-1}+(e-1)\Delta t\,\beta\bigr)\Id,
    \end{equation*}
    that is $C_{n+1}\succeq L_{n+1}\Id$. In the variable $y_n$ the recursion is
    $y_{n+1}=e^{-\Delta t}(y_n+(e-1)\Delta t)$, whose solution is the stated affine one, and
    $\Delta t\mapsto y_\infty$ is continuous and monotone on $(0,1]$ with limit $e-1$ at
    $0$ and value $1$ at $\Delta t=1$, so $y_\infty\in[1,e-1)$. The condition
    $C_n\succeq\tfrac1{2\beta}\Id$ is $y_n\leq2$, which gives $N_{\cov,-}^{\Riem}$. For the
    upper envelope, $e^X\succeq\Id+X$ again gives
    $C_{n+1}^{-1}\succeq e^{-\Delta t}(C_n^{-1}+\Delta t\,\alpha\Id)$, so
    $z_{n+1}\geq e^{-\Delta t}(z_n+\Delta t)$ and $z_n$ dominates the stated affine sequence.
    Since $z_\infty>\tfrac12$, if $z_0\geq\tfrac12$ then $z_n\geq\tfrac12$ for all $n$ as a
    convex combination of two numbers at least $\tfrac12$; if $z_0<\tfrac12$ then
    $z_0<\tfrac12<z_\infty$, the sequence increases to $z_\infty$, and $z_n\geq\tfrac12$
    first holds at the stated index. The case split is genuine, since the one-line form with
    $\logp$ is undefined for $z_0>z_\infty$.
\end{proof}

The descent estimate is proved for the retraction itself, using only a Gaussian coupling,
smoothness of the potential, Gaussian integration by parts, and a scalar exponential
inequality.

The one-step descent estimate, \Cref{lem:descent-Riem}, is stated and proved in
\Cref{app:det-proofs}.

We proceed to prove the global convergence rate of \eqref{upd:Riem}.

\begin{theorem}[Global hitting time, Riemannian scheme]\label{thm:glob-Riem}
    Under \Cref{assump:logconcave-smooth}, with
    $\lambda_{0,\min}\Id\preceq C_0\preceq\lambda_{0,\max}\Id$ and
    \begin{equation*}
        0<\Delta t\leq\frac1{2\beta\lambda_{\max}},
    \end{equation*}
    the iterates of \eqref{upd:Riem} satisfy
    $\DeltaE(a_{n+1})\leq(1-\tfrac12\alpha L_n\Delta t)\DeltaE(a_n)$ with $L_n$ the covariance
    envelope of \Cref{lem:cov-Riem}, and
    $\DeltaE(a_{n+1})\leq(1-\tfrac{\Delta t}{4\kappa})\DeltaE(a_n)$ once
    $n\geq N_{\cov,-}^{\Riem}$. In particular, for every $\delta\in(0,\DeltaE_0)$ the hitting
    time $N_{\Riem}^{\glob}(\delta):=\inf\{n\geq0:\DeltaE(a_n)\leq\delta\}$ satisfies
    \begin{equation*}
        N_{\Riem}^{\glob}(\delta)
        =
        O\!\left(
            \frac1{\Delta t}\logp\frac1{\beta\lambda_0}
            +\frac\kappa{\Delta t}\logp\frac{\DeltaE_0}\delta
        \right).
    \end{equation*}
\end{theorem}

\begin{proof}
    Apply \Cref{lem:descent-Riem} with $\Lambda=\lambda_{\max}$ and $K=\beta\lambda_{\max}$.
    Since $\Delta t\leq1/(2K)$,
    $\DeltaE(a_{n+1})\leq\DeltaE(a_n)-\tfrac{\Delta t}2\Gradnorm(a_n)^2$. If
    $C_n\succeq L_n\Id$, then \eqref{eq:FR-W2} gives
    $\Gradnorm(a_n)^2\geq\alpha L_n\DeltaE(a_n)$, and therefore
    $\DeltaE(a_{n+1})\leq(1-\tfrac12\alpha L_n\Delta t)\DeltaE(a_n)$. After
    $N_{\cov,-}^{\Riem}$ iterations \Cref{lem:cov-Riem} gives $L_n\geq1/(2\beta)$, so the
    per-step factor is at most $1-\Delta t/(4\kappa)$. The covariance burn-in costs
    $O(\Delta t^{-1}\logp(1/(\beta\lambda_0)))$ iterations, and reducing the energy gap from
    $\DeltaE_0$ to $\delta$ afterwards costs
    $O((\kappa/\Delta t)\logp(\DeltaE_0/\delta))$ iterations, using $-\log(1-x)\geq x$.
\end{proof}

\subsection{Discretization with KL divergence}

The discrete algorithm with Riemannian retraction \eqref{upd:Riem} is geometrically natural but
computationally expensive, which requires to compute a matrix exponential at every step, for
the ease in computation, we develop the following discretization scheme with KL divergence.

At the forward step, we compute
\begin{equation*}
    a_{n+1/2}
    =
    \argmin_{a\in\Aspace}
    \left\{
        \calE_2(a_n)+\inner{\nabla\calE_2(a_n)}{a-a_n}+\frac1{\Delta t}\KL(\rho_a\|\rho_{a_n})
    \right\},
\end{equation*}
which gives
\begin{equation*}
    m_{n+1/2}=m_n+\Delta t\,C_ng(a_n),
    \qquad
    C_{n+1/2}=\bigl(C_n^{-1}+\Delta t\,A(a_n)\bigr)^{-1}.
\end{equation*}
For the backward step, the update reads
\begin{equation*}
    a_{n+1}
    =
    \argmin_{a\in\Aspace}
    \left\{\calE_1(a)+\frac1{\Delta t}\KL(\rho_a\|\rho_{a_{n+1/2}})\right\},
\end{equation*}
which gives
\begin{equation*}
    m_{n+1}=m_{n+1/2},
    \qquad
    C_{n+1}=(1+\Delta t)C_{n+1/2}.
\end{equation*}

Hence, the discrete algorithm with KL divergence reads
\begin{equation}\label{upd:KL}
    \boxed{
    \begin{aligned}
        m_{n+1} &= m_n+\Delta t\,C_ng(a_n), \\
        C_{n+1} &= (1+\Delta t)\bigl(C_n^{-1}+\Delta t\,A(a_n)\bigr)^{-1},
    \end{aligned}}
\end{equation}
or equivalently $C_{n+1}^{-1}=\tfrac1{1+\Delta t}(C_n^{-1}+\Delta t\,A(a_n))$.

We denote $D(a,a_n)=\KL(\rho_a\|\rho_{a_n})$ and take note that $D(a,a_n)$ is a Bregman
divergence because
\begin{equation}\label{eq:Bregman}
    D_\psi(\xi,\xi_n)
    =\psi(\xi)-\psi(\xi_n)-\inner{\nabla\psi(\xi_n)}{\xi-\xi_n}
    =\KL(\rho_a\|\rho_{a_n}),
\end{equation}
where we define
\begin{equation}\label{eq:Bregman-potential}
    \psi(\xi)=-\frac12\log\det\bigl(\xi_2-\xi_1\xi_1^\top\bigr),
    \qquad
    \xi=(\xi_1,\xi_2)=(m,C+mm^\top).
\end{equation}
Thus \eqref{upd:KL} is a forward--backward proximal step with the Bregman penalty $D$: it
linearizes the potential part $\calE_2$ explicitly in the coordinates $(m,C)$ and treats the
entropy $\calE_1$ implicitly; only one linear solve is required per iteration. We note that
the linearization is taken in $(m,C)$, not in the expectation parameter $\xi$, so
\eqref{upd:KL} is not the exact mirror-descent step in expectation coordinates; the two maps
agree to first order in $\Delta t$. The step is the natural one for
this divergence because $\calE_2$ is relatively smooth with respect to $D$ with constant
$\beta\lambda_{\max}$, which we record as \Cref{prop:smooth-Bregman} in \Cref{app:det-proofs}.

We proceed to prove the convergence rate of the discretization algorithm with KL divergence
\eqref{upd:KL}, under the standard \Cref{assump:logconcave-smooth}. We first prove the
spectral bounds of the covariance along \eqref{upd:KL}.

\begin{lemma}[KL covariance bootstrap]\label{lem:cov-KL}
    Under \Cref{assump:logconcave-smooth}, if
    $\lambda_{0,\min}\Id\preceq C_0\preceq\lambda_{0,\max}\Id$, then along \eqref{upd:KL} one
    has $C_n\preceq\lambda_{\max}\Id$ for all $n\geq0$, and $C_n\succeq L_n\Id$ with
    \begin{equation*}
        L_0:=\min\left\{\lambda_0,\frac1{2\beta}\right\},
        \qquad
        L_{n+1}=\frac{(1+\Delta t)L_n}{1+\Delta t\,\beta L_n},
        \qquad
        L_n=\frac1{\beta\bigl[1+(\tfrac1{\beta L_0}-1)(1+\Delta t)^{-n}\bigr]} .
    \end{equation*}
    Hence $C_n\succeq\tfrac1{2\beta}\Id$ for all
    \begin{equation*}
        n\geq
        N_{\cov,-}^{\KL}
        :=\left\lceil\frac{\logp(\tfrac1{\beta L_0}-1)}{\log(1+\Delta t)}\right\rceil
        =O\!\left(\frac1{\Delta t}\logp\frac1{\beta\lambda_0}\right),
    \end{equation*}
    and $C_n\preceq\tfrac2\alpha\Id$ for all
    $n\geq N_{\cov,+}^{\KL}:=\lceil\log2/\log(1+\Delta t)\rceil=O(1/\Delta t)$.
\end{lemma}

\begin{proof}
    The update is $C_{n+1}^{-1}=\tfrac1{1+\Delta t}(C_n^{-1}+\Delta t\,A(a_n))$ with
    $\alpha\Id\preceq A(a_n)\preceq\beta\Id$. If $C_n\preceq\lambda_{\max}\Id$ then
    $C_n^{-1}\succeq\lambda_{\max}^{-1}\Id$ and
    $A(a_n)\succeq\alpha\Id\succeq\lambda_{\max}^{-1}\Id$, so
    $C_{n+1}^{-1}\succeq\lambda_{\max}^{-1}\Id$. If $C_n\succeq L_n\Id$ then
    $C_{n+1}^{-1}\preceq\tfrac1{1+\Delta t}(L_n^{-1}+\Delta t\beta)\Id$, that is
    $C_{n+1}\succeq L_{n+1}\Id$; with $y_n=(\beta L_n)^{-1}$ the recursion is
    $y_{n+1}=1+(y_n-1)/(1+\Delta t)$, so $y_n=1+(y_0-1)(1+\Delta t)^{-n}$, which is the stated
    closed form, and $y_n\leq2$ gives $N_{\cov,-}^{\KL}$. For the upper envelope,
    $A(a_n)\succeq\alpha\Id$ gives $z_{n+1}\geq1+(z_n-1)/(1+\Delta t)$ for
    $z_n:=\lambda_{\min}(C_n^{-1})/\alpha$; since the map $z\mapsto1+(z-1)/(1+\Delta t)$ is
    increasing, induction gives $z_n\geq\bar z_n:=1+(z_0-1)(1+\Delta t)^{-n}$, the comparison
    sequence solving the equality recursion. If $z_0\geq1$
    then $z_n\geq\bar z_n\geq1$ for all $n$, and if $z_0<1$ then $z_n\geq\tfrac12$ as soon as
    $(1+\Delta t)^{-n}\leq\tfrac12$, since $1-z_0\leq1$.
\end{proof}

The covariance band of \Cref{lem:cov-KL} holds for every $\Delta t>0$, in contrast with
the Riemannian bootstrap of \Cref{lem:cov-Riem}, which needs
$\Delta t\leq\min\{1,1/(\beta\lambda_{\max})\}$; the rational resolvent is unconditionally
stable in the covariance variable, and the stepsize restriction below is used only for the
descent estimate.

The forward-KL descent and Wasserstein-gradient control, \Cref{prop:onestep-KL}, is stated
and proved in \Cref{app:det-proofs}. The relevant descent divergence is the
forward divergence $\KL(\rho_{a_{n+1}}\|\rho_{a_n})$; the reverse divergence does not
satisfy the corresponding descent inequality for arbitrary mean displacement.

\begin{theorem}[Global hitting time, KL scheme]\label{thm:glob-KL}
    Under \Cref{assump:logconcave-smooth}, with
    $\lambda_{0,\min}\Id\preceq C_0\preceq\lambda_{0,\max}\Id$ and
    $0<\Delta t\leq1/(\beta\lambda_{\max})$, the iterates of \eqref{upd:KL} satisfy
    \begin{equation*}
        \DeltaE(a_{n+1})
        \leq
        \left(1-\frac{\alpha L_n\Delta t}{2(1+\Delta t\beta\lambda_{\max})^2}\right)
        \DeltaE(a_n)
    \end{equation*}
    with $L_n$ the covariance envelope of \Cref{lem:cov-KL}, and
    $\DeltaE(a_{n+1})\leq(1-\tfrac{\Delta t}{16\kappa})\DeltaE(a_n)$ once
    $n\geq N_{\cov,-}^{\KL}$. In particular, for every $\delta\in(0,\DeltaE_0)$ the hitting
    time $N_{\KL}^{\glob}(\delta):=\inf\{n\geq0:\DeltaE(a_n)\leq\delta\}$ satisfies
    \begin{equation*}
        N_{\KL}^{\glob}(\delta)
        =
        O\!\left(
            \frac1{\Delta t}\logp\frac1{\beta\lambda_0}
            +\frac\kappa{\Delta t}\logp\frac{\DeltaE_0}\delta
        \right).
    \end{equation*}
\end{theorem}

\begin{proof}
    Apply \Cref{prop:onestep-KL} with $L=L_n$, $\Lambda=\lambda_{\max}$ and
    $K=\beta\lambda_{\max}$. Combining the two displays of \Cref{prop:onestep-KL} with the
    Wasserstein Polyak--{\L}ojasiewicz inequality \eqref{eq:W2-PL} gives the stated per-step
    contraction. After $N_{\cov,-}^{\KL}$ iterations \Cref{lem:cov-KL} gives
    $L_n\geq1/(2\beta)$, and $\Delta t\beta\lambda_{\max}\leq1$ gives
    $(1+\Delta t\beta\lambda_{\max})^2\leq4$, so the per-step factor is at most
    $1-\Delta t/(16\kappa)$. The hitting-time bound follows by adding the covariance burn-in
    and using $-\log(1-x)\geq x$.
\end{proof}

The two schemes carry the same hitting-time bound and differ in cost and in stability
range. The Riemannian update requires one symmetric matrix exponential per iteration,
whereas the KL update requires only the rational resolvent
$(C_n^{-1}+\Delta t\,A(a_n))^{-1}$, that is one linear solve. The covariance band of
\eqref{upd:KL} is unconditional in $\Delta t$, while that of \eqref{upd:Riem} requires
$\Delta t\leq\min\{1,1/(\beta\lambda_{\max})\}$; the certified descent ranges are
$\Delta t\leq1/(2\beta\lambda_{\max})$ and $\Delta t\leq1/(\beta\lambda_{\max})$
respectively, so the KL scheme also admits the larger certified step.

\subsection{Sharpness of fixed-step localization}

Both hitting times carry the global-localization scale
$(\kappa/\Delta t)\log(\DeltaE_0/\delta)$, which at the globally certified stepsize
$\Delta t=\Theta(\kappa^{-1})$ becomes $\Theta(\kappa^2)$ iterations for a constant-factor
reduction of the objective. This polynomial dependence is sharp for both schemes.

\begin{theorem}[Sharp fixed-step global localization]\label{thm:sharp-disc}
    Fix $\gamma\in(0,1/2]$ and $\LH>0$. There exist constants $T,Y,q>0$ and $\kappa_0\geq1$,
    depending only on $\gamma$, $\LH$ and one fixed flat-top bump profile, such that for every
    $\kappa\geq\kappa_0$ there is an even potential $V_\kappa\in C^\infty(\R)$ with
    \begin{equation*}
        1\leq V_\kappa''\leq\kappa,
        \qquad
        \norm{V_\kappa'''}_\infty\leq\LH,
    \end{equation*}
    and an initialization
    \begin{equation*}
        m_0=\frac{Y\kappa^4}\gamma,
        \qquad
        c_0=\frac1\kappa,
    \end{equation*}
    such that, when either \eqref{upd:Riem} or \eqref{upd:KL} is run with exact expectations
    and stepsize $\Delta t=\gamma/\kappa$, one has
    \begin{equation*}
        \DeltaE(a_n)\geq q\,\DeltaE_0
        \qquad\text{for every }
        0\leq n\leq\left\lfloor\frac{T\kappa^2}\gamma\right\rfloor.
    \end{equation*}
    Consequently, for a constant-factor reduction of the objective both schemes require
    \begin{equation*}
        \Omega\!\left(\frac{\kappa^2}\gamma\right)=\Omega\!\left(\frac\kappa{\Delta t}\right)
    \end{equation*}
    iterations. The lower bound holds with $\kappastar=\kappa$, with the covariance initially
    in the curvature-scale band, and with a dimension-free Hessian-Lipschitz constant.
\end{theorem}

The construction and the proof are given in \Cref{app:bump}.

The $\Theta(\kappa^2)$ complexity reading at $\Delta t=\Theta(1/\kappa)$,
\Cref{cor:sharp-disc}, is stated in \Cref{app:bump}.

Along the shadowed trajectory $c_nA(m_n,c_n)=1+o(1)$, so both covariance maps are almost
stationary, while $c_n\simeq1/\kappa$ and $b(m_n,c_n)/m_n=\Theta(1)$. Hence the relative
mean displacement per step is
$\Delta t\,c_nb(m_n,c_n)/m_n=\Theta(\gamma/\kappa^2)$, and a constant-factor motion of the
mean requires $\Theta(\kappa^2/\gamma)$ steps. The lower bound is the product of the
globally safe fixed-step scale $\Delta t=\Theta(\kappa^{-1})$ and a Fisher--Rao global mean
scale of order $\kappa^{-1}$.

The factor $\kappa/\Delta t$ invites taking a stepsize of order one, and the bump train
alone does not forbid it: its mechanism requires the mean gain $\Delta t\,c_n$ to match the
center spacing, and a step $\kappa$ times larger jumps the mean clear of the train. The next
result closes this route. At every order-one stepsize an admissible target defeats both
schemes outright, not by slowing them down but by trapping them on an exact nonoptimal
periodic orbit.

\begin{theorem}[Exact two-cycle at order-one stepsizes]\label{thm:two-cycle}
    Fix $\LH>0$ and a stepsize
    \begin{equation*}
        \frac2\kappa<\Delta t<2\kappa .
    \end{equation*}
    Then there exist an even potential $V\in C^\infty(\R)$ with
    \begin{equation*}
        1\leq V''\leq\kappa,
        \qquad
        \norm{V'''}_\infty\leq\LH,
    \end{equation*}
    both curvature bounds attained, a mean $M>0$, and a covariance $c\in(1/\kappa,1)$ inside
    the curvature-scale band, such that both \eqref{upd:Riem} and \eqref{upd:KL}, run with
    exact expectations and stepsize $\Delta t$ from $a_0=(M,c)$, satisfy
    \begin{equation*}
        a_{2n}=(M,c),
        \qquad
        a_{2n+1}=(-M,c),
        \qquad
        n\geq0 .
    \end{equation*}
    The orbit is not optimal, $\DeltaE(a_n)=\DeltaE_{\mathrm{cyc}}>0$ for every $n$, so for
    every $\delta<\DeltaE_{\mathrm{cyc}}$ the hitting time
    $\inf\{n\geq0:\DeltaE(a_n)\leq\delta\}$ is infinite.
\end{theorem}

The construction and the proof are given in \Cref{app:cycle}. The mechanism is the
interaction of the shared explicit mean update with the covariance preconditioner. Both
covariance maps are stationary exactly where $cA(m,c)=1$, with
$A(m,c)=\E[V''(m+\sqrt cZ)]$ the averaged curvature, and there the iteration reduces to
$m^+=(1-\Delta t\,b(m,c)/(mA(m,c)))\,m$: the multiplier is governed by the ratio of the
secant slope $b(m,c)/m$, an average of the curvature over $[0,m]$, to its endpoint value
$A(m,c)$. The two are independent and both range over $[1,\kappa]$, so the ratio ranges over
$(1/\kappa,\kappa)$; prescribing it equal to $2/\Delta t$ makes the multiplier exactly $-1$,
and evenness closes the orbit. The prescription is feasible precisely when
$\Delta t>2/\kappa$, which is where \Cref{thm:sharp-disc} hands over.

\begin{corollary}[The fixed-step barrier]\label{cor:fixed-step}
    For every fixed stepsize $\Delta t>2/\kappa$ there is a target satisfying
    \Cref{assump:logconcave-smooth} with $\alpha=1$, $\beta=\kappa$ and
    $\norm{V'''}_\infty\leq\LH$, and an initialization with covariance in the curvature-scale
    band, from which neither \eqref{upd:Riem} nor \eqref{upd:KL} ever reduces the energy gap
    below a fixed positive constant. Consequently a fixed-stepsize rule that converges on
    every such target for all sufficiently large $\kappa$ must satisfy
    $\Delta t\leq2/\kappa$, where \Cref{thm:sharp-disc} forces
    $\Omega(\kappa/\Delta t)=\Omega(\kappa^2)$ iterations for a constant-factor energy
    reduction; and every rule with $\Delta t_\kappa=\Theta(1)$ fails to converge at all once
    $\kappa$ is large. The certified fixed-step scale $\Delta t=\Theta(\kappa^{-1})$ of
    \Cref{thm:glob-Riem,thm:glob-KL}, and the resulting $\Theta(\kappa^2)$ global-localization
    complexity, are therefore intrinsic to the fixed-step schemes and not artifacts of the
    proofs.
\end{corollary}

\begin{remark}[Scope, and the contrast with Bures--Wasserstein]\label{rem:cycle-scope}
    The target in \Cref{thm:two-cycle} is chosen after the stepsize, which is the quantifier
    order needed to refute a uniform fixed-step theorem; no claim is made that one target
    defeats every order-one stepsize simultaneously. Adaptivity is not ruled out in general,
    but a stability test based on the relative curvature $c_nA_n$ does not suffice: the orbit
    holds $c_nA_n=1$ exactly, so such a test accepts the step. A test based on the secant
    $b(m_n,c_n)/m_n$, a monotone line search, which rejects the flip because the energy does
    not decrease along it, and implicit or Rosenbrock-stabilized mean updates all remain
    open. The Bures--Wasserstein forward--backward scheme of
    \citet{lambert2022variational,diao2023forward} is immune for a structural reason: its
    mean update $m^+=m-\eta\,b(m,c)$ carries no covariance preconditioner, so at its
    certified stepsize $\eta\leq1/\beta$ the multiplier $1-\eta\,b/m$ lies in $[0,1)$ at
    every admissible state, uniformly in $c$, and no such orbit can form. The certified
    stepsizes of the two geometries have the same scale $\Theta(1/\kappa)$ in these units;
    the difference is that the Fisher--Rao multiplier $1-\Delta t\,c\,b/m$ degrades by the
    extra factor $c\in[1/\kappa,1]$, and it is this factor that the two-cycle exploits.
\end{remark}

\subsection{Removing the covariance burn-in}
\label{subsec:burnin}

The remaining $\logp(1/(\beta\lambda_0))$ term in
\Cref{thm:glob-cont,thm:glob-Riem,thm:glob-KL} is a genuine Fisher--Rao warm-up cost when
$C_0$ is severely underdispersed, and two optional deterministic initializers remove it. The
reason a transport step helps is the one already visible in \Cref{rem:FR-init-cov}: along the
Fisher--Rao covariance equation an eigenvalue obeys $\dot\mu_i=\mu_i(1-\mu_ih_i)$ with
$h_i\in[\alpha,\beta]$, so growth from $\lambda_0\ll1/\beta$ is multiplicative and costs time
of order $\log(1/(\beta\lambda_0))$, whereas along the Bures--Wasserstein flow
\begin{equation}\label{eq:W-flow}
    \dot m_t=g(a_t),
    \qquad
    \dot C_t=2\Id-C_tA(a_t)-A(a_t)C_t
\end{equation}
the same eigenvalue obeys $\dot\mu_i=2(1-\mu_ih_i)\approx2$, which is additive and
scale-independent.

The Bures--Wasserstein forward--backward step reads, for $a=(m,C)$ and a transport stepsize
$\eta>0$,
\begin{equation}\label{eq:BW-step}
    m^+:=m+\eta g(a),
    \qquad
    \widetilde C:=(\Id-\eta A(a))C(\Id-\eta A(a)),
    \qquad
    C^+:=\tfrac12\Bigl(\widetilde C+2\eta\Id+\bigl[\widetilde C(\widetilde C+4\eta\Id)\bigr]^{1/2}\Bigr).
\end{equation}
The covariance envelope along \eqref{eq:BW-step}, \Cref{lem:cov-BW}, is proved in
\Cref{app:det-proofs}.

\begin{theorem}[Transport-warm-started natural gradient]\label{thm:burnin-BW}
    Assume \Cref{assump:logconcave-smooth} and fix $c\in(0,1]$, with $c<1$ for the continuous
    statement, since $T_{\BW}(c)$ is finite only there.

    In continuous time, run \eqref{eq:W-flow} until $T_{\BW}(c)$ and then switch permanently to
    \eqref{ODEs:NG}. Then, for every $\delta\in(0,\DeltaE_0)$,
    \begin{equation*}
        T_{\BW\to\FR}(\delta;c)
        \leq
        T_{\BW}(c)
        +\min\left\{
            \frac\kappa c\logp\frac{\DeltaE_0}\delta,\;
            \logp\frac\kappa{c\betastar}+\kappastar\logp\frac{\DeltaE_0}\delta
        \right\}
        =
        O\!\left(\frac1\beta+\frac\kappa c\logp\frac{\DeltaE_0}\delta\right)
    \end{equation*}
    suffices for $\DeltaE(a_t)\leq\delta$.

    Discretely, set $\eta:=c/\beta$, perform one step \eqref{eq:BW-step} if
    $\lambda_0<c/\beta$ and skip it otherwise, and then run \eqref{upd:Riem} with
    $0<\Delta t\leq1/(2\beta\lambda_{\max})$, or \eqref{upd:KL} with
    $0<\Delta t\leq1/(\beta\lambda_{\max})$. Then
    \begin{equation*}
        N_{\BW\to\bullet}(\delta;c)
        =
        \mathbf1_{\{\lambda_0<c/\beta\}}
        +O\!\left(\frac\kappa{c\Delta t}\logp\frac{\DeltaE_0}\delta\right),
        \qquad \bullet\in\{\Riem,\KL\},
    \end{equation*}
    iterations suffice for $\DeltaE(a_n)\leq\delta$. In either case the covariance burn-in term
    $\logp(1/(\beta\lambda_0))$ is replaced by a single transport step, respectively by a
    warm-start time of order $1/\beta$.
\end{theorem}

\begin{proof}
    The Bures--Wasserstein stage decreases $\calE$ \citep{lambert2022variational}, so at the
    switching state $a_{\mathrm b}$ we have $\DeltaE(a_{\mathrm b})\leq\DeltaE_0$ and
    $C_{\mathrm b}\succeq(c/\beta)\Id$ by \Cref{lem:cov-BW}. In continuous time the lower
    envelope $\ell$ of \Cref{lem:cov-cont} started from $\lambda=c/\beta\leq1/\beta$ is
    nondecreasing, since
    $\ell'(t)=\lambda e^t(1-\beta\lambda)/(1+\beta\lambda(e^t-1))^2\geq0$, so
    $C_t\succeq(c/\beta)\Id$ for the whole Fisher--Rao stage; then \eqref{eq:FR-W2} and
    \eqref{eq:dissipation} give
    $\tfrac{\dd}{\dd t}\DeltaE(a_t)\leq-\tfrac c\kappa\DeltaE(a_t)$, whence the first entry in
    the minimum. Alternatively
    $\widehat C_{\mathrm b}\succeq\tfrac c\beta C_\star^{-1}\succeq\tfrac c\kappa\Id$, so
    \Cref{thm:glob-cont} restarted at $a_{\mathrm b}$ gives the second entry. The big-$O$ form
    uses $T_{\BW}(c)\leq(2\beta)^{-1}\log(1/(1-c))$ for fixed $c$.

    Discretely, \Cref{lem:cov-BW} gives $C_{\mathrm b}\succeq(c/\beta)\Id$,
    $C_{\mathrm b}\preceq\lambda_{\max}\Id$ and $\DeltaE(a_{\mathrm b})\leq\DeltaE_0$, and the
    covariance envelopes of \Cref{lem:cov-Riem,lem:cov-KL} started from
    $L_0=\min\{c/\beta,1/(2\beta)\}\geq c/(2\beta)$ are nondecreasing in $n$, since their
    $y$-variables are nonincreasing. Freezing the floor at $c/(2\beta)$ in the per-step
    contractions of \Cref{thm:glob-Riem,thm:glob-KL} gives
    $\DeltaE(a_{n+1})\leq(1-\tfrac{c\Delta t}{4\kappa})\DeltaE(a_n)$ for \eqref{upd:Riem} and
    $\DeltaE(a_{n+1})\leq(1-\tfrac{c\Delta t}{16\kappa})\DeltaE(a_n)$ for \eqref{upd:KL},
    whence the stated count.
\end{proof}

The second initializer is available whenever a pointwise Hessian can be queried at a single
point, and it lands on the full curvature band rather than only on its lower edge.

\begin{definition}[Quadratic rescue]\label{def:rescue}
    Given $m_{\mathrm{in}}\in\R^{N_\theta}$, query the oracle once at $m_{\mathrm{in}}$,
    \begin{equation*}
        G_{\mathrm{in}}:=\nabla V(m_{\mathrm{in}}),
        \qquad
        H_{\mathrm{in}}:=\nabla^2V(m_{\mathrm{in}}),
    \end{equation*}
    and set
    \begin{equation*}
        \boxed{\;
        m_0:=m_{\mathrm{in}}-H_{\mathrm{in}}^{-1}G_{\mathrm{in}},
        \qquad
        C_0:=H_{\mathrm{in}}^{-1}.\;}
    \end{equation*}
\end{definition}

The band antitonicity fact, \Cref{lem:rescue-band}, is recorded in \Cref{app:det-proofs}.

\begin{theorem}[Exact one-query recovery of Gaussian targets]\label{thm:rescue-Gauss}
    Suppose $V(\theta)=\tfrac12(\theta-\mu)^\top H(\theta-\mu)+\mathrm{const.}$ with $H\succ0$.
    Then \Cref{def:rescue} gives $m_0=\mu$ and $C_0=H^{-1}$, so $a_0=a_\star$ and
    $\DeltaE(a_0)=0$, and both \eqref{upd:Riem} and \eqref{upd:KL} remain at $a_\star$. Over
    the class of unknown Gaussian targets the oracle complexity is exactly one query.
\end{theorem}

\begin{proof}
    For a quadratic potential $\nabla V(m_{\mathrm{in}})=H(m_{\mathrm{in}}-\mu)$ and
    $\nabla^2V(m_{\mathrm{in}})=H$, so
    $m_0=m_{\mathrm{in}}-H^{-1}H(m_{\mathrm{in}}-\mu)=\mu$ and $C_0=H^{-1}$. The target belongs
    to the Gaussian family, so $a_\star=(\mu,H^{-1})=a_0$; there $g(a_\star)=0$ and
    $A(a_\star)=H=C_\star^{-1}$, so both covariance maps preserve $H^{-1}$ and both mean maps
    preserve $\mu$. For necessity, an algorithm making zero queries is independent of
    $(\mu,H)$, while one query at any $\theta$ returns $H(\theta-\mu)$ and $H$, from which
    $\mu=\theta-H^{-1}G$ and $C_\star=H^{-1}$.
\end{proof}

\begin{theorem}[Burn-in-free complexity after quadratic rescue]\label{thm:burnin-rescue}
    Assume \Cref{assump:logconcave-smooth} and initialize by \Cref{def:rescue}. Then
    $\lambda_0\geq1/\beta$, so in the covariance envelopes of \Cref{lem:cov-Riem,lem:cov-KL}
    one has $L_0=\min\{\lambda_0,\tfrac1{2\beta}\}=\tfrac1{2\beta}$ and $y_0=2$; the envelopes
    decrease from $y_0=2$ to their limits, so $C_n\succeq\tfrac1{2\beta}\Id$ for all $n\geq0$
    and $N_{\cov,-}^{\Riem}=N_{\cov,-}^{\KL}=0$. Consequently, running \eqref{upd:Riem} with
    $0<\Delta t\leq1/(2\beta\lambda_{\max})$, or \eqref{upd:KL} with
    $0<\Delta t\leq1/(\beta\lambda_{\max})$, gives for every $\delta\in(0,\DeltaE_0)$
    \begin{equation*}
        N_{\mathrm{rescue}\to\bullet}(\delta)
        =
        1+O\!\left(\frac\kappa{\Delta t}\logp\frac{\DeltaE_0}\delta\right),
        \qquad \bullet\in\{\Riem,\KL\},
    \end{equation*}
    where the additive $1$ counts the single deterministic Hessian query and
    $\DeltaE_0=\DeltaE(a_0)$ is finite, and equals $0$ for Gaussian targets by
    \Cref{thm:rescue-Gauss}. The stage of length $\logp(1/(\beta\lambda_0))$ is eliminated.
\end{theorem}

\begin{proof}
    By \Cref{lem:rescue-band}, $C_0\succeq\tfrac1\beta\Id\succ\tfrac1{2\beta}\Id$, so
    $L_0=\tfrac1{2\beta}$ and $y_0=2$. For \eqref{upd:Riem} the envelope of
    \Cref{lem:cov-Riem} is $y_n=y_\infty+(2-y_\infty)e^{-n\Delta t}$ with
    $y_\infty\in[1,e-1)\subset[1,2)$, which decreases from $2$; for \eqref{upd:KL} the envelope
    of \Cref{lem:cov-KL} is $y_n=1+(1+\Delta t)^{-n}$, which decreases from $2$ to $1$. In both
    cases $y_n\leq2$, that is $C_n\succeq\tfrac1{2\beta}\Id$, for every $n\geq0$, and the
    counts $N_{\cov,-}^\bullet=\lceil\Delta t^{-1}\logp1\rceil=0$. With the covariance floor
    $\tfrac1{2\beta}$ in force from $n=0$, the per-step contractions of
    \Cref{thm:glob-Riem,thm:glob-KL} apply from $n=0$, which gives the stated count.
\end{proof}

Both warm starts delete the $\logp(1/(\beta\lambda_0))$ stage but sit in different oracle
models and geometries. The transport warm start of \Cref{thm:burnin-BW} is
expectation-level: it uses only $g(a)$ and $A(a)$, never a pointwise Hessian, and inflates
the covariance additively, so it guarantees only the floor $C_0\succeq\tfrac c\beta\Id$ at
the lower band edge without matching the pointwise curvature. The quadratic rescue of
\Cref{def:rescue} is a single pointwise-Hessian query and lands directly on the full
curvature band $[\tfrac1\beta,\tfrac1\alpha]$, exactly recovering Gaussian targets. Thus
transport trades a weaker oracle for a coarser landing scale, while the rescue trades one
Hessian evaluation for exact Gaussian recovery; the ensuing Fisher--Rao iterations are
identical.

Both devices remove only the covariance burn-in stage. They do not improve the fixed-step
global-localization factor $\kappa/\Delta t$, since the lower construction of
\Cref{thm:sharp-disc} already starts at $c_0=1/\kappa$, after the covariance has reached
the curvature scale, and the two-cycle of \Cref{thm:two-cycle} sits at a covariance
$c\in(1/\kappa,1)$ strictly inside the band, so no covariance initialization evades either
obstruction; and they do not improve the local stiffness factor of
\Cref{sec:local}. Covariance warm-up, fixed-step global localization and local stiffness
are three distinct issues.

\subsection{Stationary guarantees for non-logconcave targets}

The convergence theory above uses \Cref{assump:logconcave-smooth} twice, once through the Wasserstein Polyak--{\L}ojasiewicz inequality on Gaussians and once through the covariance bootstrap. When $\rho_\post$ is not logconcave neither ingredient is available, but the flow \eqref{ODEs:NG} still dissipates energy at the rate of its own squared gradient norm, and this alone yields a stationarity guarantee.

\begin{corollary}\label{cor:stationary-cont}
    Suppose $\inf_{\Aspace} \calE > -\infty$. If the target distribution $\rho_\post$ is not logconcave, then along \eqref{ODEs:NG} we have the following stationary guarantee:
    \begin{equation*}
        \min_{0 \le t \le T} \| \grad^\FR \calE (a_t) \|_{a_t}^2 \le \frac{\calE (a_0) - \inf_{\Aspace} \calE}{T}.
    \end{equation*}
\end{corollary}

\begin{proof}
    Since \eqref{ODEs:NG} is a Riemannian gradient flow over the Gaussian manifold, then
    \begin{equation*}
        \frac{\dd \calE (a_t)}{\dd t} = - \| \grad^\FR \calE (a_t) \|_{a_t}^2
    \end{equation*}
    and hence
    \begin{equation*}
        \int_0^T \| \grad^\FR \calE (a_t) \|_{a_t}^2 \dd t = \calE (a_0) - \calE (a_T) \leq \calE (a_0) - \inf_{\Aspace} \calE,
    \end{equation*}
    this gives
    \begin{equation*}
        \min_{0 \le t \le T} \| \grad^\FR \calE (a_t) \|_{a_t}^2 \le \frac{\calE (a_0) - \inf_{\Aspace} \calE}{T}.
    \end{equation*}
\end{proof}

The discretizations \eqref{upd:Riem} and \eqref{upd:KL} inherit no such statement. Without \Cref{assump:logconcave-smooth} the covariance is no longer confined to the band $[L_n, \lambda_{\max}]$ of \Cref{lem:cov-Riem} and \Cref{lem:cov-KL}, and a two-sided bound on $\nabla^2 V$ alone does not restore it. We exhibit the mechanism on a smooth one-dimensional target. For $R > 0$ let
\begin{equation}\label{eq:double-well}
    V_R (\theta) := \frac{\theta^2}{2} - 2R^2 \log \cosh (\theta / R), \qquad \theta \in \R,
\end{equation}
whose derivatives are
\begin{equation*}
    V_R' (\theta) = \theta - 2R \tanh (\theta/R), \qquad V_R'' (\theta) = 1 - 2\,\mathrm{sech}^2 (\theta/R),
\end{equation*}
so that $-1 \leq V_R'' \leq 1$ uniformly in $R$. For $X \sim \N (0, c)$ we write $A_R (c) := \E [V_R'' (X)]$. Since $V_R'' = -1 + 2 \tanh^2 (\cdot / R)$ and $\tanh^2 (y) \leq y^2$, we have the quantitative curvature bound
\begin{equation}\label{eq:curvature-gap}
    0 \leq A_R (c) + 1 = 2 \E \tanh^2 (\sqrt{c} Z / R) \leq \frac{2c}{R^2}, \qquad Z \sim \N (0,1),
\end{equation}
so that $A_R (c) \to -1$ as $R \to \infty$, uniformly on compact intervals of $c$. The potential \eqref{eq:double-well} is even, hence the symmetric mean $m = 0$ is invariant under all the covariance dynamics below and only the scalar covariance moves.

\begin{proposition}[Riemannian covariance cascade]\label{prop:cascade-Riem}
    Fix $\Delta t > 0$ and let
    \begin{equation}\label{eq:cascade-map}
        c_{n+1}^{(R)} = c_n^{(R)} \exp \left\{ \Delta t \left( 1 - c_n^{(R)} A_R (c_n^{(R)}) \right) \right\}, \qquad c_0^{(R)} = 1,
    \end{equation}
    be the covariance component of \eqref{upd:Riem} for the target $\exp (-V_R)$. For every $M > 0$ there are $R$ and a finite $n$ such that $c_n^{(R)} \geq M$. Consequently no bound depending only on $\sup_\theta |V_R'' (\theta)|$, $\Delta t$ and $c_0$ can control all trajectories of \eqref{upd:Riem} uniformly over this family.
\end{proposition}

\begin{proof}
    Replacing $A_R$ by its limit $-1$ in \eqref{eq:cascade-map} gives the recursion
    \begin{equation*}
        \bar c_{n+1} = \bar c_n e^{\Delta t (1 + \bar c_n)}, \qquad \bar c_0 = 1,
    \end{equation*}
    which is increasing and diverges, so there is a finite $N$ with $\bar c_N > 2M$. By \eqref{eq:curvature-gap} the map $(R, c) \mapsto A_R (c)$ converges to $-1$ uniformly on every compact $c$-interval, and an induction over the first $N$ steps gives $c_n^{(R)} \to \bar c_n$ for $0 \leq n \leq N$ as $R \to \infty$. For $R$ sufficiently large we therefore have $c_N^{(R)} > M$.
\end{proof}

\begin{proposition}[Pole of the KL covariance resolvent]\label{prop:pole-KL}
    For $\Delta t = 1$ the covariance component of \eqref{upd:KL} for the target $\exp (-V_R)$ reads
    \begin{equation}\label{eq:KL-cov-map}
        c_1 = \frac{2 c_0}{1 + c_0 A_R (c_0)}.
    \end{equation}
    There are sequences $\varepsilon \downarrow 0$ and $R_\varepsilon \to \infty$ such that, with $c_0 = 1 - \varepsilon$,
    \begin{equation*}
        c_1 = \Theta (\varepsilon^{-1}), \qquad \left( 1 - c_1 A_{R_\varepsilon} (c_1) \right)^2 = \Theta (\varepsilon^{-2}).
    \end{equation*}
\end{proposition}

\begin{proof}
    Take $0 < \varepsilon \leq 1/4$ and choose $R_\varepsilon^2 \geq 2 \varepsilon^{-2}$. By \eqref{eq:curvature-gap}, $0 \leq A_{R_\varepsilon} (1 - \varepsilon) + 1 \leq \varepsilon^2$, so the denominator of \eqref{eq:KL-cov-map} is
    \begin{equation*}
        1 + (1 - \varepsilon) A_{R_\varepsilon} (1 - \varepsilon) = \varepsilon + (1 - \varepsilon) \left( A_{R_\varepsilon} (1 - \varepsilon) + 1 \right) = \varepsilon + O (\varepsilon^2),
    \end{equation*}
    which is positive and gives $c_1 = \Theta (\varepsilon^{-1})$, in particular $c_1 \leq 2 / \varepsilon$. Applying \eqref{eq:curvature-gap} once more at $c_1$,
    \begin{equation*}
        0 \leq A_{R_\varepsilon} (c_1) + 1 \leq \frac{2 c_1}{R_\varepsilon^2} \leq 2 \varepsilon \leq \frac{1}{2},
    \end{equation*}
    whence $1 + c_1 / 2 \leq 1 - c_1 A_{R_\varepsilon} (c_1) = 1 + c_1 - c_1 (A_{R_\varepsilon} (c_1) + 1) \leq 1 + c_1$, and the residual scaling follows.
\end{proof}

The blow-up in \eqref{eq:KL-cov-map} is caused by the approaching loss of positive definiteness of the rational resolvent, and it occurs despite the uniform Hessian bound $|V_R''| \leq 1$; taken together, the two propositions show that a global smoothness constant cannot replace covariance control in a trajectory-wise Fisher--Rao theorem.

\begin{remark}
    The Bures--Wasserstein forward--backward method carries a different guarantee beyond log-concavity. Under the assumptions and stepsize $\eta$ of \citet[Theorem~5.5]{diao2023forward}, its iterates satisfy the running-minimum bound
    \begin{equation*}
        \min_{0 \leq n < N} \| \grad^\BW \calE (a_n) \|_\BW^2 \leq \frac{150 \{ \calE (a_0) - \inf \calE \}}{\eta N}.
    \end{equation*}
    This is a stationarity statement in the Bures--Wasserstein geometry and does not control the pointwise Fisher--Rao residual appearing in \Cref{prop:cascade-Riem,prop:pole-KL}. We use it only as a comparison between the two geometries.
\end{remark}

\subsection{A covariance-constrained KL scheme}

\Cref{prop:cascade-Riem,prop:pole-KL} suggest constraining the covariance directly, as in \citet{sun2025natural}. Fix constants $0 < \lambda_- \leq \lambda_+ < \infty$ and define the covariance-truncated Gaussian parameter space
\begin{equation*}
    \Aspace' := \left\{ a = (m, C) \in \R^{N_\theta} \times \SPD (N_\theta) : \lambda_- I \preceq C \preceq \lambda_+ I \right\},
\end{equation*}
and assume the following global smoothness condition throughout this subsection.

\begin{assumption}\label{assump:smooth}
    The potential $V \in C^2 (\R^{N_\theta})$ satisfies
    \begin{equation*}
        \| \nabla^2 V (\theta) \|_\op \leq \beta, \qquad \forall \theta \in \R^{N_\theta},
    \end{equation*}
    equivalently $-\beta I \preceq \nabla^2 V (\theta) \preceq \beta I$.
\end{assumption}

For $C \in \SPD (N_\theta)$ with spectral decomposition $C = Q \diag (c_1, \ldots, c_{N_\theta}) Q^\top$ we define the clipping operator
\begin{equation*}
    \Clip_{[\lambda_-, \lambda_+]} (C) := Q \diag \left( \min \{ \lambda_+, \max \{ \lambda_-, c_1 \} \}, \ldots, \min \{ \lambda_+, \max \{ \lambda_-, c_{N_\theta} \} \} \right) Q^\top.
\end{equation*}
Clipping the spectrum is exactly the Bregman projection onto the covariance interval for the log-determinant divergence $D_\Phi (\cdot, \cdot)$ generated by $\Phi (C) = -\frac{1}{2} \log \det C$, which is the mechanism that lets the projected step retain a descent certificate.

At $a_n = (m_n, C_n)$ we write
\begin{equation*}
    b_n := \E_{\N (m_n, C_n)} [\nabla_\theta V (\theta)], \qquad A_n := \E_{\N (m_n, C_n)} [\nabla_\theta \nabla_\theta V (\theta)],
\end{equation*}
so that \eqref{upd:KL} followed by the projection reads
\begin{equation}\label{upd:clip-KL}
    \begin{aligned}
        m_{n+1} &= m_n - \Delta t\, C_n b_n, \\
        \widetilde{C}_{n+1} &= (1 + \Delta t) (C_n^{-1} + \Delta t A_n)^{-1}, \qquad C_{n+1} = \Clip_{[\lambda_-, \lambda_+]} (\widetilde{C}_{n+1}).
    \end{aligned}
\end{equation}
The inverse is well defined at the stepsizes considered below. The one-step stationarity measure of this split scheme is
\begin{equation}\label{eq:split-stationarity}
    \mathsf{S}_n := \frac{1}{2} \| m_{n+1} - m_n \|_{C_n^{-1}}^2 + \KL \left( \N (0, C_n) \, \middle\| \, \N (0, C_{n+1}) \right),
\end{equation}
where the orientation of the covariance divergence matters: the forward divergence, with $C_{n+1}$ in the first slot, carries the smoothness remainder, while the reverse one appearing in \eqref{eq:split-stationarity} is what the descent certificate leaves behind.

\begin{theorem}\label{thm:clip-KL}
    Under \Cref{assump:smooth}, if $a_0 \in \Aspace'$ and the stepsize satisfies $0 < \Delta t \leq \frac{1}{2 \beta \lambda_+}$, then the iterates of \eqref{upd:clip-KL} are well defined, remain in $\Aspace'$, and satisfy
    \begin{equation}\label{eq:clip-descent}
        \calE (a_{n+1}) \leq \calE (a_n) - \frac{1}{\Delta t} \mathsf{S}_n.
    \end{equation}
    Consequently, for every $N \geq 1$,
    \begin{equation}\label{eq:clip-runmin}
        \min_{0 \leq n < N} \mathsf{S}_n \leq \frac{\Delta t}{N} \left\{ \calE (a_0) - \calE (a_N) \right\}.
    \end{equation}
\end{theorem}

\begin{proof}
    Since $\lambda_- I \preceq C_n \preceq \lambda_+ I$ and $\| A_n \|_\op \leq \beta$, the stepsize restriction gives
    \begin{equation*}
        C_n^{-1} + \Delta t A_n \succeq (\lambda_+^{-1} - \Delta t \beta) I \succeq (2 \lambda_+)^{-1} I,
    \end{equation*}
    so $\widetilde{C}_{n+1}$ is well defined, and $C_{n+1} \in [\lambda_- I, \lambda_+ I]$ by construction. Writing $\calE = \Phi + \calE_2 + \mathrm{const}$ with $\Phi (C) = -\frac{1}{2} \log \det C$ and $\calE_2 (a) = \E_{\rho_a} [V (\theta)]$, a square-root coupling of $\rho_{a_n}$ and $\rho_{a_{n+1}}$ together with the scalar inequality $x - 1 - \log x \geq (\sqrt{x} - 1)^2$ shows that $\calE_2$ is relatively smooth with constant $L = 2 \beta \lambda_+$ with respect to the split divergence built from $\frac12 \| \cdot \|_{C_n^{-1}}^2$ and $\KL (\N (0, \cdot) \| \N (0, C_n))$. The pair $(m_{n+1}, C_{n+1})$ is the unique minimizer over $\Aspace'$ of the associated proximal model, because the covariance part of that model is, up to an additive constant, a multiple of $D_\Phi (\cdot, \widetilde{C}_{n+1})$, whose Bregman projection onto $[\lambda_- I, \lambda_+ I]$ is spectral clipping. Mean optimality then contributes $b_n^\top (m_{n+1} - m_n) = -\frac{2}{\Delta t} \cdot \frac12 \| m_{n+1} - m_n \|_{C_n^{-1}}^2$, while the constrained log-determinant three-point inequality contributes the reverse covariance divergence in \eqref{eq:split-stationarity}. Since $L \Delta t \leq 1$, the forward remainder is absorbed and the two estimates combine to \eqref{eq:clip-descent}. Summing \eqref{eq:clip-descent} over $0 \leq n < N$ and bounding the minimum by the average gives \eqref{eq:clip-runmin}. The details are given in \Cref{app:nonconvex}.
\end{proof}

\begin{corollary}\label{cor:clip-KL-full}
    Under the assumptions of \Cref{thm:clip-KL}, the full successive-Gaussian reverse KL divergence satisfies
    \begin{equation*}
        \min_{0 \leq n < N} \KL \left( \N (m_n, C_n) \, \middle\| \, \N (m_{n+1}, C_{n+1}) \right) \leq \frac{\lambda_+}{\lambda_-} \frac{\Delta t}{N} \left\{ \calE (a_0) - \calE (a_N) \right\}.
    \end{equation*}
\end{corollary}

\begin{proof}
    The reverse KL divergence between successive Gaussians splits as
    \begin{equation*}
        \KL \left( \N (m_n, C_n) \, \middle\| \, \N (m_{n+1}, C_{n+1}) \right) = \KL \left( \N (0, C_n) \, \middle\| \, \N (0, C_{n+1}) \right) + \frac{1}{2} \| m_{n+1} - m_n \|_{C_{n+1}^{-1}}^2,
    \end{equation*}
    and the mean block is measured in $C_{n+1}^{-1}$ rather than $C_n^{-1}$. Feasibility gives
    \begin{equation*}
        C_{n+1}^{-1} \preceq \lambda_-^{-1} I \preceq \frac{\lambda_+}{\lambda_-} C_n^{-1},
    \end{equation*}
    so the full divergence is at most $(\lambda_+ / \lambda_-) \mathsf{S}_n$, and \eqref{eq:clip-runmin} completes the proof.
\end{proof}

For centered scalar trajectories the mean term vanishes, so the factor-one bound \eqref{eq:clip-runmin} is the quantity plotted in \Cref{sec:numerics-nonconvex}.

\begin{remark}
    \Cref{thm:clip-KL} certifies stationarity of the projected split map on the covariance-truncated set $\Aspace'$. If an iterate is pinned at the upper clip $\lambda_+ I$, then $\mathsf{S}_n$ can vanish while the unconstrained Fisher--Rao gradient points outward. The theorem therefore does not imply that the unconstrained gradient is small, nor that the iterates converge to the unconstrained Gaussian optimizer $a_\star$.
\end{remark}

% SOURCES:
% natural-gradient.tex l.198-220 (Corollary, continuous stationarity) -> cor:stationary-cont (verbatim modulo notation/labels)
% codex-draft/sections/06-nonconvex.tex l.13-36 (eq:nonconvex-potential, A_R, limit) -> eq:double-well, A_R
% codex-draft/sections/06-nonconvex.tex l.92-98 (eq:nonconvex-curvature-quantitative) -> eq:curvature-gap (hoisted out of the pole proof, used by both propositions)
% codex-draft/sections/06-nonconvex.tex l.39-67 (prop:nonconvex-riemannian-cascade + proof) -> prop:cascade-Riem, eq:cascade-map
% codex-draft/sections/06-nonconvex.tex l.69-119 (prop:nonconvex-kl-pole + proof) -> prop:pole-KL, eq:KL-cov-map
% codex-draft/sections/06-nonconvex.tex l.121-125 (closing moral) -> single closing sentence
% codex-draft/sections/06-nonconvex.tex l.127-145 (eq:nonconvex-bw-running-min, Diao Thm 5.5) -> BW comparison remark
% natural-gradient.tex l.1048-1073 (Aspace', assump:smooth, Clip) -> Aspace', assump:smooth, Clip
% codex-draft/sections/06-nonconvex.tex l.150-201 (split update, split stationarity measure) -> upd:clip-KL, eq:split-stationarity
% codex-draft/sections/06-nonconvex.tex l.203-251 (thm:nonconvex-clipped-stationarity + proof) -> thm:clip-KL
%   (supersedes natural-gradient.tex l.1095-1122, whose proof does not handle the clipping/three-point interaction)
% codex-draft/sections/06-nonconvex.tex l.253-286 (cor:nonconvex-full-kl + proof) -> cor:clip-KL-full
%   (its closing sentence -> pointer to \Cref{sec:numerics-nonconvex} in sections/06-numerics.tex)
% codex-draft/sections/06-nonconvex.tex l.288-296 (rem:nonconvex-scope) -> scope remark
%   (the stochastic clipped theorem of natural-gradient.tex l.1137-1195 is deliberately not carried)


% SOURCES:
%   §2.1
%     natural-gradient.tex l.111-124 (Assumption, section opening prose) -> assump:logconcave-smooth
%     improved-global-local-stoch.tex l.2489-2494 (assump:global-Hess-Lip, copied verbatim;
%       moved ahead of its first use per the restructure) -> assump:hess-lip
%     improved-global-local.tex l.339-395 (g(a), A(a), R(a), first-order conditions, splitting) -> eq:gA, eq:first-order, eq:E-split
%     improved-global-local.tex l.494-616 (lem:metric-dissipation) -> eq:FR-metric, eq:differential, eq:FR-gradient, eq:grad-norm, eq:dissipation
%     natural-gradient.tex l.329-332 (FR metric display, voice) -> eq:FR-metric
%     improved-global-local.tex l.620-691 (lem:W2-PL, lem:FR-W2) -> eq:W2-PL, eq:FR-W2
%       [these supersede the red TODO tcolorbox at natural-gradient.tex l.155, which is dropped]
%   §2.2
%     improved-global-local.tex l.821-843 (lem:precision-duhamel) -> eq:precision-duhamel
%     improved-global-local.tex l.845-883 (lem:cont-cov) -> lem:cov-cont
%     improved-global-local.tex l.700-817 (lem:optimizer-whitening, lem:kappa-star-kappa) -> eq:whitening, eq:star-curvature, eq:lamzstar, eq:kappa-star-kappa
%     improved-global-local.tex l.885-962 (thm:cont-global) -> thm:glob-cont
%       [supersedes natural-gradient.tex thm:glob-conv, which is not carried]
%     improved-global-local.tex l.964-971 (rem: what kappa_star encodes) -> remark after thm:glob-cont
%     natural-gradient.tex l.222-305 (rem:FR-init-cov, kept verbatim modulo notation;
%       final sentence redirected from the cut WFR section to §2.6) -> rem:FR-init-cov
%     improved-global-local.tex l.1614-1655 (thm:spiral-sharp) + l.1550-1612 (lem:spiral-kappa-star)
%       + l.1236-1262 (lem:spiral-shadowing radius condition) -> thm:sharp-cont
%     improved-global-local.tex l.1657-1683 (cor:cont-sharp) -> cor:sharp-cont
%     improved-global-local.tex l.1685-1715 (Hessian-Lipschitz and scope remarks) -> scope remark
%   §2.3
%     natural-gradient.tex l.308-362 (splitting, base-point penalty, forward/backward steps,
%       boxed update; kept near-verbatim) -> eq:Riem-dist, upd:Riem
%       [DEFECT FIXED: the source calls eq:Riem-dist "the distance of the metric" and uses an
%        unqualified Log map; the blockwise penalty is not the Fisher-Rao geodesic distance
%        (mean block frozen at C_n; not symmetric in its arguments). Relabeled here as a
%        base-point penalty and the Log replaced by the inverse retraction Retr^{-1},
%        consistent with rem:Riem-retr and improved-global-local.tex l.1845-1847.]
%     natural-gradient.tex l.366-397 (prop:upd-map-Riem) -> prop:upd-map-Riem
%     improved-global-local.tex l.1845-1847 ("exponential retraction") -> rem:Riem-retr
%     improved-global-local.tex l.1719-1839 (lem:Riem-cov) -> lem:cov-Riem
%     improved-global-local.tex l.1879-1991 (lem:Riem-descent) -> lem:descent-Riem (proof in app:det-proofs)
%     improved-global-local.tex l.1993-2037 (thm:Riem-global) -> thm:glob-Riem
%       [supersedes natural-gradient.tex thm:conv-Riem, which is not carried]
%   §2.4
%     natural-gradient.tex l.504-541 (motivating sentence, forward/backward steps, boxed update,
%       Bregman identification with psi; kept near-verbatim) -> upd:KL, eq:Bregman, eq:Bregman-potential
%     improved-global-local-stoch.tex l.3994-4025 (lem:KL-argmin) -> forward-backward proximal sentence
%       [DEFECT FIXED: the source calls upd:KL "a forward-backward mirror step in the
%        expectation parameters", but its linearization of E_2 is in (m,C), not in
%        xi=(m,C+mm^T); the exact xi-mirror step is a different map (agrees only to first
%        order in Dt). Reworded as a Bregman-penalty proximal step with the discrepancy
%        stated explicitly.]
%     improved-global-local.tex l.2041-2104 (lem:KL-cov) -> lem:cov-KL
%     improved-global-local.tex l.2164-2285 (prop:KL-onestep) -> prop:onestep-KL (proof in app:det-proofs)
%     improved-global-local.tex l.2110-2113 (forward vs reverse divergence) -> sentence after prop:onestep-KL
%     improved-global-local.tex l.2287-2325 (thm:KL-global) -> thm:glob-KL
%       [supersedes natural-gradient.tex thm:conv-KL, which is not carried]
%   §2.5
%     improved-global-local.tex l.2918-2969 (thm:disc-global-sharp) -> thm:sharp-disc
%     improved-global-local.tex l.2971-2983 (cor:disc-global-sharp) -> cor:sharp-disc
%     improved-global-local.tex l.2985-2999 (rem: why the construction is slow) -> mechanism remark
%     improved-global-local.tex l.3001-3008 (rem: limitations) -> scope remark
%     improved-global-local.tex l.2330-2342 (localization scale, lead-in) -> lead-in sentence
%   §2.6
%     improved-global-local.tex l.4315-4351 / improved-global-local-stoch.tex l.1818-1829
%       (multiplicative vs additive covariance growth) -> lead-in contrast sentence
%     improved-global-local.tex l.4382-4452 (lem:W-cov-bootstrap) + l.4587-4642 (lem:one-step-W-bootstrap)
%       -> lem:cov-BW, eq:W-flow, eq:BW-step
%     improved-global-local.tex l.4454-4554 (thm:cont-W-to-FR, global stage only)
%       + l.4644-4751 (thm:disc-W-to-FR, global stage only) -> thm:burnin-BW
%     improved-global-local-stoch.tex l.1908-1931 (def:rescue, lem:rescue-band) -> def:rescue, lem:rescue-band
%     improved-global-local-stoch.tex l.1933-1948 (thm:gauss-recovery) -> thm:rescue-Gauss
%     improved-global-local-stoch.tex l.1950-2004 (thm:disc-rescue-to-FR, global stage only)
%       + l.2231-2234 (lem:rescue-energy, finiteness of Delta_0) -> thm:burnin-rescue
%     improved-global-local-stoch.tex l.2006-2018 (rem:two-initializers) -> oracle-tradeoff remark
%     improved-global-local.tex l.4753-4764 (rem:transport-bootstrap-caveat) -> caveat remark
